\section{Classical and Romantic Ethics}

\begin{quotex}
Before I formed thee in the bowels of thy mother, I knew thee: and before thou camest forth out of the womb, I sanctified thee, and made thee a prophet unto the nations. \flright{\textsc{Jeremiah 1:5}}

\end{quotex}
In the chapter Classical and Romantic Ethics from \textit{Sintesi di dottrina della razza}, \textbf{Julius Evola} recapitulates the traditional doctrine of preexistence. This was one of the acceptable opinions for the origin of the soul along with traducianism and creationism. The latter developed in opposition of the former, but has nothing to say against preexistence, properly understood (i.e., not in a crude, ``materialistic" sense). Every possibility of manifestation ``exists" in the mind of God. If you want to restrict the word ``exist" to manifested being, then that possibility does not ``exist" as an actuality. Moreover, if by ``preexistence" is meant a prior existence as another human being, then that use of the word is false. In a planned post on the metaphysics of \textbf{Vladimir Solovyov}, we will defend the use of the neoplatonic and other traditional sources cited by Evola.

Evola himself refers to the Catholic teaching on the creation of the soul from nothing, with the same semantic confusion about ``preexistence" mentioned above. This chapter cannot be understood in the conventional scientific, empirical, or material sense, but requires a full understanding of the metaphysical principles underlying it. Thus ``race" in his view is not at all a biological quality that can be determined genetically, but a spiritual attitude. Otherwise the decline of the modern world would be incomprehensible. It results from a spiritual decline, regardless of biological facts.

These are questions that are ignored by theologians today. The idea that archangels are assigned to different peoples and nations is no longer taken seriously, never mind that they participate in their creation. The quote from Plotinus on inequality used to be known and was held to be due to God's Providence, predestination, and predilection. Of course, our intent is not to open a debate on these topics, but rather to observe their effects. All things, including men, have their origin in the mind of God as possibilities. If a man has a destiny, he must actively actualize it and not wait passively for it to happen. In \textbf{Rene Guenon}`s words, that refers to the actualization of the possibilities open to him. That is, he must become who he is.


\hfill

\paragraph{Classical and Romantic Ethics}
In the work just cited we reproduced various traditional texts that clarify and confirm those views. In particular, in this regard, there is this passage from \textbf{Plotinus}:

\begin{quotex}
The general plane is unique, but is divided into different parts, so that in all there are distinct places, some more pleasant and others less—and souls, also unequal, come to inhabit these distinct places that correspond to their own differences. In such a way everything is concordant and the difference of the situations corresponds to the inequality of the souls. 

\end{quotex}
And even more precisely:

\begin{quotex}
The soul chose itself prior to its daemon and its life. 

\end{quotex}
\textbf{Plato} similarly taught:

\begin{quotex}
It is not the daemon that chooses you, but you yourself are to choose the daemon. You yourself choose the destiny of that life to which you then will be irremediably connected. 

\end{quotex}
These last expressions are particularly interesting for us, since the concept of daemon has nothing to do with the Christian concept of an evil entity, but has instead the closest relation with the deepest forces of race, both of the soul and of the body. We cannot here analyze the traditional doctrine in this regard, but only recall that in the same way the ``daemon", the ``\emph{lares}", the ``\emph{penates}", the ``double" (which is a synonym of the ``subtle body") are notions that in antiquity merged and reflected the precise knowledge of the true roots of the differentiations of bloods, of \emph{gentes}, and ultimately of individuals themselves, on the basis of a complete vision of the world, comprising the invisible and the visible, and not of that modern distortion which knows only about material and psychological processes. From such evidence, which could be multiplied with reference to the traditions of all peoples, we therefore see confirmed the idea of transcendental, or vertical, heredity and of the choice that, on the basis of analogical correspondences, determines its connection to ``horizontal", historical-biological heredity. The consequences of all that in regard to the justification of the idea of race theory are quite obvious.

The central view of Catholicism is that God, while creating man out of nothing, let the miracle happen, through which this being created from nothing is free, in the sense that he can rejoin God at the root of his own being, or else deny Him, be limited to himself, disperse himself, or degenerate into a useless creaturely free will. This same doctrine, with due transpositions, can be applied to the relations between the individual being and the spiritual entity, of which it is the creation and human manifestation. We mean that the individual being, within given limits, enjoys equally free will and that the same alternative is posed to him: either to will his own nature, deepening it and realizing it to the point of rejoining himself to the pre-human and superindividual principle that corresponds to him; or to give himself up to arbitrarily construct an unnatural mode of being, deprived of relation with one's deepest forces or directly in contradiction with them. This is exactly the existing opposition between the traditional and especially the Nordic-Aryan ideal and the ``modern" ideal of civilization. For the former, the essential task is to know and be oneself; for the latter, the task is instead to ``construct oneself", to become what one is not, to break every limit to render possible everything to everyone: liberalism, democracy, individualism, activistic-protestant ethics, anti-race, anti-traditionalism.

As was traditionally taught, the doctrine of preexistence therefore leads beyond both fatalism, as well as a poorly understood and individualistic freedom. Moving on to more immediate consequences, in realizing his own nature, the individual harmonizes his own human will with the super-human will which corresponds to him, he ``remembers" himself, he establishes the relation with a principle that, being beyond birth, is likewise beyond death and every temporal condition: therefore, according to the ancient Indo-Aryan conception, this is the way through which he aims to achieve ``liberation" and to realize the divine, by way of action. Dharma—which means one's own nature, duty, fidelity to blood, tradition, caste—is connected, as we already explained in the other book, to the sensation of having arrived here from far away and does not mean limitation, as the ``evolved spirits" believe, but liberation. Led back to this traditional vision of life, all the principle themes of race theory acquire a higher and spiritual significance and the objection based on birth as chance or destiny loses every force.

But that is not enough: it is not by chance that the formula, ``know yourself",  that in its deepest meaning precisely refers to such teachings, and was itself written on Apollo's Delphic temple, i.e., of the Hyperborean God. To let such traditional truths act on oneself, to the point where they awaken precise inner forces, means to proceed on the way that leads to a spiritual level from which the meaning of life constitutes something absolutely different from that of the rest of mankind: the importance of clarity, absolute strength, and incomparable certainty. But to have a presentiment of all that, to barely see a ``style" in which the feeling of detachment of ``those who have reached from afar" and of interior inaccessibility unites to himself a type of indomitability, in which there is simultaneously a superior serenity, a distance, and a readiness to attack, to command, to absolute action—to have a presentiment of this ``style", means also to have concealed the mystery of the primordial Nordic, or Hyperborean race, as race of the spirit. Such is in fact the Olympic and solar way of being. Popular imagination relates it today to the so called ``men of destiny" and previously related it to the rare types of great leaders—in reality, in what were the last echoes or flashes of what was typical, in general, of the great Hyperborean super-race, before its dispersion and degradation. We recall Plutarch's expression about the same members of the ancient Roman Senate: ``They sit like a council of kings".

An additional consequence follows from this: if a civilization of the ``classical" type, in this Olympic and virile sense, not in the vulgar esthetic and formalistic meaning, reflects something of the Nordic race of the spirit, every romantic and tragic civilization, which is opposed to it, will be instead the certain sign of influences taking priority, that proceeded from racial and ethnic residues of a non-Nordic, pre-Aryan, and anti-Aryan nature.



\flrightit{Posted on 2015-07-16 by Cologero }

\begin{center}* * *\end{center}

\begin{footnotesize}\begin{sffamily}



\texttt{Mark Citadel on 2015-07-18 at 13:10 said: }

Man's pre-existence in the mind of God as a possibility (with this constituting a true existence) would seem to rule out all arguments for abortion, and then some. Those not yet born are most certainly the same as those whom they will become, with their passage of birth being a mere semantic.

The Christian view of abortion is affirmed entirely here.


\hfill

\texttt{Cologero on 2015-07-18 at 16:39 said: }

Agreed, Mr Citadel. This doctrine is stronger and, to my mind, preferable to the ``creationist" account which seems rather arbitrary and ad hoc.


\hfill

\texttt{Matt on 2015-07-19 at 15:55 said: }

It's certainly a different view of pre-existence when compared to what others have written on the subject, with the exception being Guenon, I guess.

From my readings of Evola on the subject of pre-existence, I always got the impression that his view of the pre-existence of the particular subject was the same as that of a Origen, or a Saint-Martin (Martin's being more along the lines of the Kabbalah's Adam-Kadmon) . That the soul/person pre-exists embodiment as a manifested pure spirit, not just as a possibility of manifestation in the Divine Mind.

More material to contemplate on.


\hfill


\hfill

\texttt{Cologero on 2015-07-19 at 20:56 said: }

Matt, I took the liberty of bringing out Evola's real meaning while putting it into a larger context. As a literal event, Evola's notion of a soul choosing its incarnation does not make much sense. Nevertheless, his rethought notions of karma and dharma are powerful. So we can understand the created soul as having qualities; this is also consistent with predilection. Moreover, the categories of being, at least at the moment of conception, are not simply ``accidents", but are essential to the soul's being.


\end{sffamily}\end{footnotesize}
